\subsection{Learned VAE priors and explicit program relationships}
\label{sec:vaepriors}
Learned latent priors are an established component of variational
autoencoders. Their relevance here is both statistical and computational:
they provide alternatives for representing dependence, but differ in the
meaning assigned to latent coordinates and in the inference operations
available under the observation model. Three distinctions are essential:
prior flexibility does not imply an interpretable dependency graph;
conditional independence within a mixture component does not imply
marginal independence; and a tractable prior density does not imply a
tractable posterior under an arbitrary decoder.

\paragraph{VampPrior: learned pseudo-inputs and shared mixture components.}
Tomczak and Welling~\cite{vampprior} define the basic VampPrior by
\begin{equation}\label{eq:vampposition}
p(z)=\frac1H\sum_{h=1}^Hq_\phi(z\mid u_h),
\end{equation}
where \(q_\phi\) is the VAE encoder and the \(u_h\) are trainable vectors
in the observation space (their equation~9). Pseudo-inputs are global
model parameters optimized with the encoder and decoder. They need not
be realistic observations, and they are not noisy copies of fitted
loadings or sample-specific latent variables.
With a diagonal Gaussian encoder, \eqref{eq:vampposition} is a vector
Gaussian mixture whose parameters are tied through the encoder. It is
not a normalizing flow.

Equivalently, introduce a sample-specific label
\(C_i\sim\operatorname{Uniform}\{1,\ldots,H\}\) and
\(z_i\mid C_i=h\sim q_\phi(z_i\mid u_h)\).
The label is shared across coordinates of \(z_i\), so the mixture can
encode nonlinear and multimodal conditional dependence even when each
component has diagonal covariance. Section~\ref{sec:twineb} derives the
corresponding coordinate conditionals. Neither diagonal components nor
a diagonal variational encoder justify describing this population prior
as independent. The latter restricts the approximate posterior for each
observation, not the dependence of the mixture prior.

The same paper separately introduces a two-layer hierarchical VAE, with a
VampPrior on its upper latent layer and a neural Gaussian conditional
distribution for the lower layer (their equations~12--15).
That extra continuous latent layer is not required by the basic VampPrior.
In contrast, the present hierarchy specifies dependencies among the
existing factorization loadings. Its mixture indicators are auxiliary
labels for the specified component distributions, rather than a new
continuous representation above the programs. The HMM states serve the
separate feature-sequence prior.

\paragraph{Autoregressive-flow priors.}
The Variational Lossy Autoencoder of Chen et al.~\cite{vlae} parameterizes
a latent prior by an invertible autoregressive transformation
\(z=f_\theta(\epsilon)\) of a simple continuous noise source.
Its density follows from the change-of-variables identity,
\begin{equation}\label{eq:flowposition}
\log p_\theta(z)=
\log p_0(f_\theta^{-1}(z))
+\log\left|\det J_{f_\theta^{-1}}(z)\right|.
\end{equation}
Such a prior explicitly introduces dependence between latent coordinates.
A single affine Gaussian autoregressive step has Gaussian scalar
conditionals in its ordering, but this restriction should not be
generalized to arbitrary compositions and nonlinear flow families.
Accordingly, complex distributional dependence itself is not a
distinguishing capability of the present model.

The distinction is the specified semantics of its conditional heads:
an exact activation probability, a chosen amplitude distribution, and
an explicit parent set for each additive program. A smooth invertible
flow with an absolutely continuous base distribution assigns no positive
mass to an exact zero event. Exact spikes therefore require an additional
discrete or mixed-support construction; an ordinary continuous flow
does not provide them by itself. Conversely, adding a neural conditional
mixture does not automatically give more interpretable effects than a
flow. Interpretation requires restrictions on the heads and stable
definitions of the latent programs.

\paragraph{iVAE and conditional identifiability.}
Khemakhem et al.~\cite{ivae} study priors of the form
\begin{equation}\label{eq:ivaeposition}
p(z\mid x)=\prod_k p_k(z_k\mid x),
\end{equation}
with suitable exponential-family component distributions and an observed
auxiliary variable \(x\). Conditional factorization, together with
additional assumptions on the decoder and variation of the conditional
parameters, supports their identifiability results. It does not assert
marginal independence after averaging over \(x\).
Their use of \emph{exponential family} refers to the general statistical
class, not specifically to the exponential slabs in
Section~\ref{sec:mixturefamilies}.

The present model allows direct dependence on latent parents even after
conditioning on \(x\). It therefore addresses a different modeling
objective and does not automatically inherit iVAE's identification
guarantees. A directed loading prior is not a causal identification
result, and coefficient interpretation remains conditional on the chosen
ordering, scales and program definitions.

\paragraph{An interpretable linear-logistic specialization.}
Let \(Z^R_{ik}=1\) indicate the continuous slab of RNA loading \(R_{ik}\);
when \(Z^R_{ik}=0\), set \(R_{ik}=0\). A restricted CGB prior can specify
\begin{equation}\label{eq:interpretablegate}
\logit\Pr(Z^R_{ik}=1\mid A_i,R_{i,<k},x_i)
=\alpha_k+\beta_k^\top A_i
             +\eta_k^\top R_{i,<k}+\gamma_k^\top x_i .
\end{equation}
For a parent-independent Gaussian slab, the coefficients govern activation
probability rather than slab amplitude. Holding all other parents and
context fixed, increasing \(A_{ij}\) by \(\Delta\) multiplies the prior
activation odds by \(\exp(\beta_{kj}\Delta)\).
This is a conditional association at a specified loading scale, not a
causal effect or the posterior odds after observing both modalities.
The latter also include the likelihood and child contributions derived
earlier. Sign-constrained Gaussian means do not, by themselves, make
loadings positive.

Equation~\eqref{eq:interpretablegate} gives a more direct coefficient
interpretation than a generic vector-mixture or flow parameterization.
It does not prove superior empirical interpretability: latent rotations,
permutations, scale changes and alternative program decompositions can
change the apparent relationships. Anchoring or aligning programs and
assessing stability across fits remain necessary for scientific use.
The parent set consists of the prescribed preceding nodes. Unrestricted
regressions of every loading on all the other loadings are not
automatically compatible with a single joint distribution.

The current experimental CGB implementation uses ReLU neural gates.
Equation~\eqref{eq:interpretablegate} is a restricted model specification,
not a coefficient interpretation of those fitted neural networks.
Section~\ref{sec:linearlogistic} describes its simpler, log-concave slab
update under the stated assumptions. Neural gates and EMDN can subsequently
be compared with this restriction when nonlinear activation, amplitude or
variance relationships are scientifically required.

\paragraph{Observed context in a vector-mixture alternative.}
A conditional population mixture can incorporate the same fixed context:
\begin{align}
C_i\mid x_i&\sim\operatorname{Categorical}(w_1(x_i),\ldots,w_H(x_i)),\\
v_i\mid C_i=h,x_i&\sim\N(\mu_h(x_i),D_h(x_i)),
\qquad v_i=(A_i^\top,R_i^\top)^\top .
\label{eq:contextvectormixture}
\end{align}
Context may affect mixture probabilities, component means or covariance
matrices; it need not become an uncertain latent variable.
A conditional VampPrior could tie components through an encoder
\(q_\phi(v\mid u_h,x)\), whereas a directly parameterized mixture can
learn these component functions without introducing that encoder.
These are possible extensions, not a claim that the original VampPrior
paper implements this multi-omics model.

Under diagonal component covariances, dependence between programs is
mediated by the shared component label. Their conditional relationships
can still be complex, but an effect of one program on another is usually
distributed across the component parameters and responsibilities rather
than summarized by a coefficient such as \(\beta_{kj}\).
Full component covariances permit additional dependence within each
component.

\paragraph{Why Gaussian mixtures can simplify inference.}
Conditional on the feature factors, the observed row vector satisfies
a linear Gaussian model \(Y_i^{\rm obs}=B_i v_i+\epsilon_i\), where
\(B_i\) selects the observed entries from the two modality designs and
\(\epsilon_i\sim\N(0,\Sigma_i)\).
For positive-definite Gaussian prior covariances in
\eqref{eq:contextvectormixture}, the row conditional is a Gaussian mixture.
Its component parameters are
\begin{align}
S_{ih}&=(D_h^{-1}+B_i^\top\Sigma_i^{-1}B_i)^{-1},\\
m_{ih}&=S_{ih}(D_h^{-1}\mu_h+B_i^\top\Sigma_i^{-1}Y_i^{\rm obs}),
\end{align}
with weights proportional to
\[
w_h(x_i)\,
\N\!\left(Y_i^{\rm obs};B_i\mu_h,
               \Sigma_i+B_iD_hB_i^\top\right).
\]
Here \(\mu_h,D_h\) are evaluated at fixed \(x_i\).
These identities follow by completing the multivariate Gaussian square.
They establish tractability of a row conditional at fixed factors and
parameters, not an analytic posterior for the entire factorization.
Missing entries or modalities are omitted from the observation vector
and design. HMM factor updates remain available conditional on the
loading draws.

Both this mixture model and the direct hierarchy propagate information
between modalities. The computational simplification is Gaussian
conjugacy of the joint components; it does not arise merely from adding
a latent label. In the direct logistic hierarchy, an upstream update
instead contains nonlinear log-gate contributions from its children.
Those contributions preserve the explicit model's semantics while
making its update nonconjugate. Exact sparsity patterns in vector
mixtures require explicit lower-dimensional component supports, and
exponential components do not generally preserve the unrestricted
Gaussian block formulas above. Neither alternative is uniformly cheaper
at comparable modeling accuracy.

\paragraph{Positioning and empirical consequences.}
For an interpretation-oriented study, the proposed main model is a
direct conditional loading hierarchy, with linear-logistic activation
heads as a transparent restriction. Its contribution must be assessed
through interpretable conditional associations, useful predictions and
joint uncertainty propagation, rather than through a claim that previous
VAE priors cannot encode dependence.
Conditional vector mixtures offer a relevant tractable comparator,
and neural gates offer an expressiveness comparison within the same
factorization framework. Differences in observation models, feature
priors, ranks and tuning budgets must be controlled.

Finally, replacing the linear factorization by a nonlinear VAE decoder
preserves the normalized conditional-prior model and its ELBO, but
generally removes the exact quadratic normal-means likelihood.
The analytic updates in this paper are therefore not automatically
available for arbitrary VAEs. Learned priors and structured posterior
approximations are related choices, but neither posterior dependence
nor exact spike probabilities follow from changing a prior alone.
